\sectioncentered*{Введение}
\addcontentsline{toc}{section}{Введение}

Большие языковые модели сделали возможным управление программными системами на
естественном языке, что особенно органично для интерактивных виртуальных сред.
Однако у ассистентов, построенных исключительно на языковой модели, обнаруживается
принципиальное ограничение: модель хранит знания в весах, полученных при обучении,
и не располагает ни актуальным состоянием среды, ни механизмом проверки собственных
утверждений. Как следствие, она порождает правдоподобные, но неверные
идентификаторы объектов и не способна гарантировать соблюдение ограничений. Ошибка
при этом совершается молча: внешне ответ неотличим от корректного, и пользователь
не получает сигнала о том, что его поняли неправильно\cite{Ji2023}.

Разрешение противоречия достигается объединением двух подходов в единой архитектуре,
при котором логико-символьный подход отвечает за хранение знаний, вывод и объяснение,
а нейробионический~--~за понимание свободных формулировок. Системы такого рода
называются гибридными и рассматриваются как одно из основных
направлений развития искусственного интеллекта. Актуальность темы определяется тем,
что вынесение знаний и проверок за пределы языковой модели даёт гарантии корректности,
недостижимые статистическими методами, при сохранении естественности взаимодействия.
Задача возникает всюду, где языковая модель управляет действиями в системе с
собственным состоянием и собственными правилами\cite{Garcez2023}.

Целью курсовой работы является разработка гибридной интеллектуальной системы
игрового ассистента, объединяющей онтологическую базу знаний, решатель задач и
большую языковую модель в едином процессе решения задач пользователя.

Для достижения цели необходимо решить следующие задачи:

\begin{itemize}
  \item проанализировать предметную область, класс решаемых задач и существующие
        гибридные интеллектуальные системы;
  \item определить состав и назначение компонентов системы: базы знаний, решателя
        задач и интерфейса;
  \item сформировать требования к гибридной интеллектуальной системе;
  \item спроектировать архитектуру, определить точки интеграции компонентов и
        стратегию их совместного применения;
  \item реализовать базу знаний, решатель задач и интеграцию компонентов;
  \item провести тестирование и сравнительную оценку гибридного, чисто символьного
        и чисто нейросетевого вариантов.
\end{itemize}

Логико-символьный подход представлен онтологической базой знаний в формате RDF/OWL,
правилами ограничений и решателем, выполняющим обратный вывод по графу способов
получения объектов; нейробионический~--~большой языковой моделью, которая извлекает
намерение из свободной реплики, подбирает официальные наименования для разговорных и
формулирует ответ по выходу решателя.

Объектом исследования являются гибридные интеллектуальные системы, объединяющие
логико-символьные и нейросетевые методы. Предметом исследования являются методы
интеграции онтологической базы знаний и решателя задач с большой языковой моделью
в диалоговой системе управления действиями.

Практическая значимость работы состоит в том, что предложенная архитектура
обеспечивает нулевую долю недопустимых действий по построению, а не по статистике:
решение о допустимости принимается решателем после языковой модели и не может быть
изменено формулировкой запроса.
